\input{../tex-inputs/latex-leseansicht-vorspann}

\section[Gustav Schwarzkopf an Arthur Schnitzler, 29. 7. 1899]{L04297 Gustav Schwarzkopf an Arthur Schnitzler, 29. 7. 1899}
\nopagebreak\mylabel{L04297v}
\rehead{ }\normalsize\beginnumbering\briefempfaengerindex{Schnitzler, Arthur@\textsc{Schnitzler, Arthur}!zzzSchwarzkopf, Gustav@\emph{von Gustav Schwarzkopf}!1899-07-292@{29. 7. 1899}|(be}
\toendnotes[C]{\smallbreak\pagebreak[2]}
\correspDesc{Versand  durch Gustav Schwarzkopf am 29. 7. 1899 in Wien
\newline{}Erhalt  durch Arthur Schnitzler im Zeitraum [31. 7. 1899 – 3. 8. 1899?] in Toblach}\toendnotes[C]{\smallbreak}
\Standort{CUL, Schnitzler, B 96a.}
\physDesc{Brief, 1 Blatt, 4 Seiten, 3068 Zeichen
\newline{}Handschrift: schwarze Tinte, deutsche Kurrent}\toendnotes[C]{\smallbreak}
\pstart
           \raggedleft{}{\pb}29. 7. 99\pend
           \vspace{0.5em}
\pstart
           Lieber Arthur, ich war erſtaunt noch aus \label{K_L04297-1v}\edtext{\textsc{Velden\oindex{Velden am Wörthersee@\textbf{Velden am Wörthersee}|pw}} einen Brief}{\lemma{\textnormal{\emph{Velden einen Brief}}}\Cendnote{\textnormal{XXXX Auszeichnungsfehler: Dokument L04043 nicht gefunden.}}}\label{K_L04297-1} zu beko{\geminationm}en, ich hab geglaubt, daß Sie{ }ſchon in \textsc{Ischl\oindex{Bad Ischl@\textbf{Bad Ischl}|pw}}{ }ſind. D\textsuperscript{r.}\textsc{Kapper\pwindex{Kapper, Hans 10.\,3.\,1891 Wien – September 1962 New York City@\textsc{Kapper, Hans} (10.\,3.\,1891 Wien – September 1962 New York City), \emph{Schriftsteller}|pw}}{ }ſagte mir, daß Sie nach \textsc{Ischl\oindex{Bad Ischl@\textbf{Bad Ischl}|pw}} berufen würden zu ihrer \label{K_L04297-2v}\edtext{Patientin\pwindex{Schey, Marie 8.\,5.\,1821 Nagykanizsa – 22.\,8.\,1899 Bad Ischl@\textsc{Schey, Marie} (8.\,5.\,1821 Nagykanizsa – 22.\,8.\,1899 Bad Ischl)|pwv}}{\lemma{\textnormal{\emph{Patientin}}}\Cendnote{\textnormal{Schnitzlers Tante Marie Schey\pwindex{Schey, Marie 8.\,5.\,1821 Nagykanizsa – 22.\,8.\,1899 Bad Ischl@\textsc{Schey, Marie} (8.\,5.\,1821 Nagykanizsa – 22.\,8.\,1899 Bad Ischl)|pwk} war
                  seit 10. 3. 1898 seine Patientin. Sie starb am 
                  22. 8. 1899 in Ischl\oindex{Bad Ischl@\textbf{Bad Ischl}|pwk}.}}}\label{K_L04297-2}, der alten Dame\pwindex{Schey, Marie 8.\,5.\,1821 Nagykanizsa – 22.\,8.\,1899 Bad Ischl@\textsc{Schey, Marie} (8.\,5.\,1821 Nagykanizsa – 22.\,8.\,1899 Bad Ischl)|pwv}
               aus der Praterſtraße\oindex{Wien@\textbf{Wien}!II., Leopoldstadt@\textbf{II., Leopoldstadt}!Praterstraße 54@\textbf{Praterstraße 54}, \emph{Wohngebäude}|pwv}, die gefährlich erkrankt{ }ſei. Seine Frau\pwindex{Schnitzler, Olga 17.\,1.\,1882 Wien – 13.\,1.\,1970 Lugano@\textsc{Schnitzler, Olga} (17.\,1.\,1882 Wien – 13.\,1.\,1970 Lugano), \emph{Schauspielerin, Sängerin}|pwv} hat es ihm
               geſchrieben, die es von Ihrem Cousin Mandl\pwindex{Mandl, Ludwig 19.\,8.\,1862 Iași – 22.\,9.\,1937 Wien@\textsc{Mandl, Ludwig} (19.\,8.\,1862 Iași – 22.\,9.\,1937 Wien), \emph{Gynäkologe}|pw}
               erfahren hatte. Es war alſo nicht{ }ſo{ }ſchli{\geminationm} oder
               überhaupt nur Gerücht?\pend
           
\pstart
           Es{ }ſcheint, daß Sie{ }ſich in \textsc{Velden\oindex{Velden am Wörthersee@\textbf{Velden am Wörthersee}|pw}} doch ziemlich behaglich gefühlt haben. Der Ort muß{ }ſich wol{ }ſehr zu{ }ſeinem
               Vorteil verändert haben. Damals hat man{ }ſo{ }ſchlecht gegeſſen. Geſtern hab ich Hofmansthal’s\pwindex{Hofmannsthal, Hugo August von 21.\,12.\,1841 Wien – 8.\,12.\,1915 ebd.@\textsc{Hofmannsthal, Hugo August von} (21.\,12.\,1841 Wien – 8.\,12.\,1915 ebd.), \emph{Bankdirektor}|pw}\pwindex{Hofmannsthal, Anna von 27.\,1.\,1849 Wien – 22.\,3.\,1904 Sanatorium Fürth@\textsc{Hofmannsthal, Anna von} (27.\,1.\,1849 Wien – 22.\,3.\,1904 Sanatorium Fürth)|pw}\pwindex{Hofmannsthal, Hugo von 1.\,2.\,1874 Wien – 15.\,7.\,1929 Rodaun@\textsc{Hofmannsthal, Hugo von} (1.\,2.\,1874 Wien – 15.\,7.\,1929 Rodaun), \emph{Schriftsteller}|pw} geſprochen,
               die {\pb}mit ihrem Aufenthalt in Marienbad\oindex{Marienbad@\textbf{Marienbad}|pw} { }ſehr unzufrieden waren und \strikeout{ihn} vor der beſti{\geminationm}ten Zeit
               fort{ }ſind. Hier auf der Ringſtraße\oindex{Wien@\textbf{Wien}!I., Innere Stadt@\textbf{I., Innere Stadt}!Ringstraße@\textbf{Ringstraße}, \emph{Straße}|pw}, der noch i{\geminationm}er meine Abendſpaziergänge gelten, iſt man wenigſtens
               vor Enttäuſchungen{ }ſicher. Der liebe, gute, alte Staub bleibt{ }ſich dort i{\geminationm}er gleich. Auf einer dieſer So{\geminationm}erreiſen habe ich auch D\textsuperscript{r}Kapper\pwindex{Kapper, Hans 10.\,3.\,1891 Wien – September 1962 New York City@\textsc{Kapper, Hans} (10.\,3.\,1891 Wien – September 1962 New York City), \emph{Schriftsteller}|pw} getroffen, geſtern auch \textsc{Eberma{\geminationn}\pwindex{Ebermann, Leo 16.\,7.\,1863 Draganovka – 9.\,10.\,1914 Wien@\textsc{Ebermann, Leo} (16.\,7.\,1863 Draganovka – 9.\,10.\,1914 Wien), \emph{Schriftsteller, Journalist, Rechtswissenschaftler}|pw}} und \textsc{Beraton\pwindex{Bératon, Ferry 6.\,12.\,1859 Wien – 11.\,2.\,1900 Venedig@\textsc{Bératon, Ferry} (6.\,12.\,1859 Wien – 11.\,2.\,1900 Venedig), \emph{Schriftsteller, Journalist, Maler}|pw}}. Oh, es{ }ſind{ }ſchon noch Leute in Wien\oindex{Wien@\textbf{Wien}, \emph{Verwaltungsgebiet}|pw},{ }ſogar
               \label{K_L04297-3v}\edtext{Ihr \textsc{Hirschfeld\pwindex{Hirschfeld, Georg 11.\,2.\,1873 Berlin – 17.\,1.\,1942 München@\textsc{Hirschfeld, Georg} (11.\,2.\,1873 Berlin – 17.\,1.\,1942 München), \emph{Schriftsteller}|pw}}}{\lemma{\textnormal{\emph{Ihr Hirschfeld}}}\Cendnote{\textnormal{Vermutlich meint er Georg Hirschfeld\pwindex{Hirschfeld, Georg 11.\,2.\,1873 Berlin – 17.\,1.\,1942 München@\textsc{Hirschfeld, Georg} (11.\,2.\,1873 Berlin – 17.\,1.\,1942 München), \emph{Schriftsteller}|pwk}, wohingegen 
                  Robert Hirschfeld\pwindex{Hirschfeld, Robert 17.\,9.\,1857 Žďár nad Sázavou – 2.\,4.\,1914 Salzburg@\textsc{Hirschfeld, Robert} (17.\,9.\,1857 Žďár nad Sázavou – 2.\,4.\,1914 Salzburg), \emph{Journalist, Musikkritiker}|pwk} eher als Schwarzkopfs\pwindex{Schwarzkopf, Gustav 7.\,11.\,1853 Wien – 13.\,11.\,1939 ebd.@\textsc{Schwarzkopf, Gustav} (7.\,11.\,1853 Wien – 13.\,11.\,1939 ebd.), \emph{Schriftsteller}|pwk} Freund
                  gegolten haben dürfte.
               }}}\label{K_L04297-3}. –\pend
           
\pstart
           Übrigens bin ich auch einmal mit der neuen Stadtbahn\orgindex{Wiener Stadtbahn@Wiener Stadtbahn|pw} von Karlsplatz\oindex{Karlsplatz@\textbf{Karlsplatz}, \emph{Platz}|pw} nach Hietzing\oindex{III., Landstraße@\textbf{III., Landstraße}, \emph{Verwaltungsgebiet}|pw} gefahren und hab dabei an unſere letzte
               »grausige Fahrt« gedacht.\pend
           
\pstart
           Auch für Abwechslung iſt geſorgt, Max\pwindex{Schwarzkopf, Max 12.\,6.\,1857 Wien – 14.\,4.\,1928 ebd.@\textsc{Schwarzkopf, Max} (12.\,6.\,1857 Wien – 14.\,4.\,1928 ebd.), \emph{Rechtsanwalt}|pw} iſt
               heute angeko{\geminationm}en (ſehr zufrieden) Emil\pwindex{Schwarzkopf, Emil 17.\,9.\,1851 Wien – 28.\,1.\,1928 ebd.@\textsc{Schwarzkopf, Emil} (17.\,9.\,1851 Wien – 28.\,1.\,1928 ebd.), \emph{Übersetzer, Komponist}|pw} fährt übermorgen für 6 Wochen fort.\pend
           
\pstart
           Was mich betrifft,{ }ſo glaube ich nicht, daß ich die berühmte Fußpartie, {\pb}die wol zum großen Teil im Landauer
               zurückgelegt werden dürfte, mitmachen werde. Wenn ich 4 Seiten mit den Gründen
               anfüllen wollte, die mich davon abhalten,{ }ſo würde das nicht viel nützen, da Sie
               dieſe Gründe doch unzureichend finden würden;{ }ſie sind es vielleicht auch wirklich,
               ich erſpare{ }ſie alſo Ihnen und mir. Es wäre ja ganz leicht, ich kö{\geminationn}te einfach{ }ſagen, daß ich gerade vom 6.
               bis ausgerechnet 16. Auguſt von einem heftigen Arbeitsfieber befallen{ }ſein werde, das mir nicht geſtattet u.ſ.w., aber das würden Sie mir nicht glauben;{ }ſagen wir alſo, es iſt Sti{\geminationm}ungsſache. (Ich weiß, daß ich
               eigentlich gar kein Recht habe, Sti{\geminationm}ungen unterworfen zu{ }ſein, daß man dazu viel jünger{ }ſein und noch{ }ſo und{ }ſoviele andere Eigenſchaften
               haben muß, aber manchmals lebe ich halt auch über meine Verhältniſſe.)\pend
           
\pstart
           Im Vertrauen teile ich Ihnen mit, daß ich mich mit der kühnen Idee trage {\pb}gar nicht fortzugehen, wenigſtens
               nicht für längere Zeit, auf die Geſahr hin mich um allen Credit zu bringen und im
               nächſten Herbſt, we{\geminationn} das Geſpräch mit den Leiden und
               Freuden des So{\geminationm}ers beſtritten wird, ganz
               gesellſchaftsunfähig zu werden. Es fehlt mir die Luſt und Energie mich der
               Einförmigkeit, an die ich mich bereit\textcolor{gray}{s} gewöhnt habe, zu entreißen,
               die ich da{\geminationn} nach einer kurzen Unterbrechung wieder{ }ſtärker empfinden würde, an die ich mich neu gewöhnen müßte. Vielleicht ko{\geminationm} ich aber doch auf zwei – drei Tage nach \textsc{Ischl\oindex{Bad Ischl@\textbf{Bad Ischl}|pw}}; für dieſe Auszeichung würde{ }ſich aber \textsc{Ischl\oindex{Bad Ischl@\textbf{Bad Ischl}|pw}} da{\geminationn} nur bei Ihnen zu bedanken haben. Konſtatiren
               will ich auch noch, daß dies der längſte Brief iſt, den ich seit Jahren geſchrieben
               habe. Nun wünſche ich Ihnen und Ihren Reiſegefährten\pwindex{Beer-Hofmann, Richard 11.\,7.\,1866 Wien – 26.\,9.\,1945 New York City@\textsc{Beer-Hofmann, Richard} (11.\,7.\,1866 Wien – 26.\,9.\,1945 New York City), \emph{Schriftsteller}|pwv}\pwindex{Wassermann, Jakob 10.\,3.\,1873 Fürth – 1.\,1.\,1934 Altaussee@\textsc{Wassermann, Jakob} (10.\,3.\,1873 Fürth – 1.\,1.\,1934 Altaussee), \emph{Schriftsteller}|pwv}\pwindex{Hirschfeld, Robert 17.\,9.\,1857 Žďár nad Sázavou – 2.\,4.\,1914 Salzburg@\textsc{Hirschfeld, Robert} (17.\,9.\,1857 Žďár nad Sázavou – 2.\,4.\,1914 Salzburg), \emph{Journalist, Musikkritiker}|pwv}, denen ich meine ergebenſten Grüße zu melden bitte,{ }ſchönes
               Wetter und viel Vergnügen.\pend
           
\pstart
           Herzlichſt{\\[\baselineskip]}\spacefill\mbox{Gustav.}\pend
           \leftskip=0em{}\selectlanguage{ngerman}\endnumbering\briefempfaengerindex{Schnitzler, Arthur@\textsc{Schnitzler, Arthur}!zzzSchwarzkopf, Gustav@\emph{von Gustav Schwarzkopf}!1899-07-292@{29. 7. 1899}|)be}\mylabel{L04297h}
\begin{anhang}
\end{anhang}\newcommand{\dateiname}{L04297}\newcommand{\titel}{Gustav Schwarzkopf an Arthur Schnitzler, 29. 7. 1899}\newcommand{\editorInnen}{Herausgegeben von Jahnke, SelmaMüller, Martin Anton}\input{../tex-inputs/latex-leseansicht-abspann}
